%!TEX root=../main.tex

\aaron{TODO: remove this section and incorporate text into ReLU section.}

Let \( X \in \R^{n \times d} \) be a data matrix, \( y \in \R^n \) the
associated vector of targets,
and \( \calB = \cbr{b_1, \ldots, b_m} \) a disjoint partition of the feature indices
\( \cbr{1, \ldots, d} \).
Given a regularization parameter \( \lambda \geq 0 \), the linear group lasso
solves the following regularized regression problem:
\begin{equation}\label{eq:group-lasso}
	p^*(\lambda) = \min_{w} \half \norm{X w - y}_2^2 + \lambda \sum_{b_i \in \calB} \norm{\wi}_2.
\end{equation}
Solutions to \cref{eq:group-lasso} are block sparse when \( \lambda \) is sufficiently large, meaning \( \wi = 0 \) for some subset of \( \bi \).
This is similar to the feature sparsity given by the lasso,
to which the group lasso naturally reduces when \( \bi = \cbr{i} \) for each
\( \bi \in \calB \).

Our primary interest in this work is the solution map or
``regularization path'',
\[
	\solfn(\lambda) := \argmin_w f_\lambda(w).
\]
For a general data matrix, \( f_\lambda \) is not strictly convex and
the linear group lasso problem may admit multiple solutions.
As such, \( \solfn \) is a point-to-set map and
only the min-norm solution,
\[
	\wmin(\lambda) = \argmin \cbr{ \norm{w}_2 : w \in \solfn(\lambda) },
\]
defines a function.
Even in the non-unique case, the model fit \( \hat y(\lambda) = X w \) is the
same for all \( w \in \solfn(\lambda) \) \citep[Lemma 2]{vaiter2012dof}.
Moreover, if \( \lambda > 0 \), then it is straightforward to deduce that the sum of group norms
\( \sum_{i \in \calB} \norm{\wi}_2 \) is also constant over \( \solfn(\lambda) \).
Uniqueness of \( \hat y(\lambda) \) extends to the residual vector
\( r(\lambda) = y - \hat y(\lambda) \)
and the block correlation vectors \( \ci(\lambda) = \Xbi^\top r(\lambda) \).
We write \( c \in \R^d \) as the concatenation of the \( \ci \)
blocks.

Unlike the lasso, the group lasso solution map is not piece-wise linear
in general \citep{yuan2006model}.
This prevents us from obtaining an explicit expression for the solution
path between transition point, where features enter or exit the active set.
We overcome this difficulty by relying only on an \emph{implicit} characterization
of the path provided by the implicit function theorem.
We start towards this goal by developing optimality conditions for the group
lasso problem.

Since \cref{eq:group-lasso} is a convex optimization problem, first-order (FO)
optimality conditions are both necessary and sufficient for weights \( w \)
to be globally optimal.
Sub-differentiating the optimization problem gives the following FO conditions:
\begin{equation}\label{eq:fo-optimality}
	\begin{aligned}
		\ci(\lambda) := \Xbi^\top (y - X w) \in
		\begin{cases}
			\cbr{\lambda \frac{\wi}{\norm{\wi}_2}} & \mbox{if \( \wi \neq 0 \)} \\
			\cbr{v : \norm{v}_2 \leq \lambda }     & \mbox{otherwise,}
		\end{cases}
	\end{aligned}
\end{equation}
showing that \( \wi = 0 \) when \( \norm{\ci}_2 < \lambda \).
As a result, the equicorrelation set
\begin{equation}\label{eq:equi}
	\equi := \cbr{\bi \in \calB : \norm{\Xbi^\top (y - Xw)}_2 = \lambda},
\end{equation}
contains all blocks which may be active for fixed \( \lambda \).
The complement of \( \equi \) in \( \calB \) is the inactive set,
denoted \( \inact \).
Our first result combines FO conditions with uniqueness of \( \hat y(\lambda) \)
to characterize the solution map in terms of \( \Null(\Xe) \).
\begin{restatable}{proposition}{solutionFunctionAlt}\label{prop:sol-fn-alt}
	Suppose \( \lambda > 0 \).
	Then the solution function for the group lasso is given by
	\begin{equation}
		\begin{aligned}
			\solfn(\lambda) =
			\bigg\{ w \in \R^d : \,
			 & \forall \, \bi \in \equi, \, \wi = \alpha_\bi \ci, \alpha_\bi \geq 0, \, \\
			 & \forall \, b_j \in \calB \setminus \equi, \w_{b_j} = 0, \,
			X w = \hat y.
			\bigg\}
		\end{aligned}
	\end{equation}
\end{restatable}
As an immediate consequence, the solution map is a particular subset
of directions in \( \Null(\Xe) \).
\begin{corollary}\label{cor:uniqueness-cond}
	If \( w, w' \in \solfn \),
	then \cref{prop:sol-fn-alt} implies
	\[
		w - w' \in \calN_\lambda := \Null(\Xe) \cap \cbr{z : \zi \in \Span(\ci), i \in \equi }.
	\]
	As a result, the group lasso solution is unique when \( \calN_\lambda = \emptyset \).
\end{corollary}
\cref{cor:uniqueness-cond} extends a similar result for the lasso
to the group lasso \citep[Eq. 9]{tibshirani2013unique} and
implies the group lasso solution is unique for all \( \lambda \geq 0 \) if
the columns of \( X \) are linearly independent;
however, we can also use it to obtain the following refined condition
for uniqueness.
%\begin{restatable}{proposition}{GGP}[Group General Position]\label{prop:unique-sol}
%    Suppose for every \( \calE \subseteq \calB \), \( \abs{\calE} \leq n +1 \),
%    there do not exist unit vectors \( \zi \in \R^{\abs{\bi}} \)
%    such that for any \( j \in \calE \),
%    \[
%        X_{b_j} z_{b_j} \in
%        \text{affine}(\cbr{\Xbi \zi : \bi \in \calE \setminus b_j}).
%    \]
%    Then the group lasso solution is unique for all \( \lambda > 0 \).
%\end{restatable}
We call this uniqueness condition \emph{group general position}
(GGP) because it naturally extends general position to groups of column vectors.
General position itself is an extension of affine independence and is sufficient for the lasso solution
to be unique \citep{tibshirani2013unique}.
We note that GGP is strictly weaker than linear independence of the columns of \( X \),
but neither implies nor is implied by general position (\cref{prop:general-position-relation}).

\subsection{Properties of the Solution Map}

Now we prove several general results about the solution map \( \solfn \).
First we provide standard definitions for continuity-type properties of
a point-to-set map.

\begin{definition}[Closed Map]
	A point-to-set map \( T : \calX \into \calZ \) is closed if
	\( \cbr{\xk} \subset \calX \), \( \xk \into \bar x \) and \( z_k \in T(\xk) \),
	\( z_k \into \bar z \) implies \( \bar z \in T(\bar x) \).
\end{definition}

\begin{definition}[Open Map]
	A point-to-set map \( T : \calX \into \calZ \) is open if
	\( \cbr{\xk} \subset \calX \), \( \xk \into \bar x \) and \( \bar z \in T(\bar x) \),
	implies there exists \( k' \in \bbN \), \( z_k \in T(\xk) \) for \( k \geq k' \),
	such that \( z_k \into \bar z \).
\end{definition}

We say that \( T \) is \emph{continuous} if it is both closed and open.
If \( T(x) \) is a singleton for all \( x \in \calX \),
then openness/closedness of \( T \) are equivalent to the standard
notion of continuity.

Our first result of this section establishes continuity of the optimal
objective value.

\begin{restatable}{proposition}{valueContinuity}\label{prop:value-continuity}
	The optimal function value \( p^* \) is continuous for all
	\( \lambda \geq 0 \).
\end{restatable}

By analyzing the dual problem, we can extend continuity to the
model fit and the sum of group norms.

%\begin{restatable}{proposition}{modelFitContinuity}\label{prop:model-fit-continuity}
%    The optimal model fit \( \hat y(\lambda)  \) is continuous for all \( \lambda \geq 0 \) the penalty \( \sum_{\bi \in \calB} \norm{\wi(\lambda)}_2 \) is continuous
%    for \( \lambda > 0 \).
%\end{restatable}

It is straightforward to extend this characterization to closedness
of \( \solfn \), but openness is not possible in the general setting.

\begin{restatable}{proposition}{mapContinuity}\label{prop:map-continuity}
	The solution map \( \solfn \) is closed on \( \R_+ \)
	and open on \( R^+ \) if only if \( X \) has full column rank.
	However, if GGP holds, then \( \solfn \) is open for every \( \lambda > 0 \).
\end{restatable}

Continuity of \( \solfn \) is impossible in the general setting because, as \citet{hogan1973point} shows, openness is essentially a stability property, one which cannot be satisfies when the equicorrelation set changes.
Continuity of the unique solution path is an immediate corollary of
these developments.

\begin{corollary}
	The group lasso solution is unique and continuous for every \( \lambda > 0 \)
	when GGP holds and is also continuous at \( \lambda = 0 \) when \( X \)
	is full column rank.
\end{corollary}

